\section{Phương pháp thực hiện}
Quy trình thực hiện đề tài được chia thành năm bước: khảo sát bài toán, chọn method, thiết kế kiến trúc, triển khai thử nghiệm và đánh giá định tính. Hình \ref{fig:research-process} mô tả quy trình thực hiện ở mức tổng quát.

\begin{figure}[h]
\centering
\begin{tikzpicture}[
  node distance=1.05cm,
  processstep/.style={rectangle,draw=Navy,rounded corners,fill=LightBlue,align=center,minimum width=3.0cm,minimum height=0.85cm},
  arrow/.style={-{Latex[length=2mm]},thick,Navy}
]
\node[processstep] (s1) {1. Khảo sát\\bài toán};
\node[processstep,right=of s1] (s2) {2. Chọn\\method};
\node[processstep,right=of s2] (s3) {3. Thiết kế\\hệ thống};
\node[processstep,below=of s3] (s4) {4. Triển khai\\thử nghiệm};
\node[processstep,left=of s4] (s5) {5. Đánh giá\\và cải tiến};
\draw[arrow] (s1) -- (s2);
\draw[arrow] (s2) -- (s3);
\draw[arrow] (s3) -- (s4);
\draw[arrow] (s4) -- (s5);
\end{tikzpicture}
\caption{Quy trình thực hiện đề tài SourceVerify}
\label{fig:research-process}
\end{figure}
